% =====================================================================
%  EurekaZero -- Section 5: The Price of Play
%  File: sections/05_price.tex   (compiles under the shared preamble)
% =====================================================================
\section{The Price of Play}\label{sec:price}
% Reader question: Can anyone afford to play it?
% One-sentence answer: Yes; cost per win is cost per attempt divided by the
%   win rate, so wasted moves set the bill, not token prices.
% Core points advanced: 3 (search = 1) and 2 (search length is the
%   reciprocal of taste); supports 1 (cheap replay is what invention lacks).
% Exhibits: Equation 2 (eq:bill), Table 5 (tab:price), Figure 4 (fig:price), Box 3.
% Every number that is not cited comes from scripts/cost_model.py.

Can anyone afford to play the innovation game? Yes: we project that a full
training run with a 30B-class open-weight program writer costs about \$12,000
in GPU time, and that wasted moves, not token prices, set the bill. Cheap
replay is the one property of Go that invention lacks (\cref{sec:game}). Here
we price it, and \cref{sec:results} sets a measurement beside each projection.

\subsection{The bill in one line}\label{sub:price-bill}

A win costs what an attempt costs, repairs included, divided by the chance
that an attempt wins:
\begin{equation}
  C_{\text{win}} = \frac{C_{\text{attempt}} + \mathbb{E}[r]\,C_{\text{repair}}}{q}.
  \label{eq:bill}
\end{equation}
Here $C_{\text{attempt}}$ is the cost of one try, including the referee's
fit; $C_{\text{repair}}$ is the cost of one repair turn, and $\mathbb{E}[r]$
the mean number of repairs per attempt. The $q$ is that of \cref{eq:taste},
the probability that one try beats the incumbent: what sets the number of
tries also sets the price. It factors as $q = q_{\text{run}}\,q_{\text{beat}}$,
the share of moves that run times the share of those that win. A move that
runs and loses still costs money, so we divide by $q$, not by the survival
rate. Generating a token costs about 2~FLOPs per active parameter, and
training on one about~6 \citep{kaplan2020scaling}. We price every run both
self-hosted and at API list prices, each dated (\cref{sub:price-run}).

The bill has two ledgers. Training pays \cref{eq:bill} millions of times;
game day pays for one attempt and one fit, with no repair, whatever $q$ is. Agent swarms buy
judgment by the try on every new problem, while self-play pays for it once,
and \cref{sec:results} measures where the two ledgers cross.

\subsection{Wasted moves}\label{sub:price-waste}

Most moves are wasted, in two ways. A stillborn move never competes: it fails
to parse or type-check, crashes, times out or returns degenerate output. An
also-ran runs and loses. \Cref{fig:price}a follows pilot moves through four
gates and records the compute spent on the moves that stop at each. At 90\%
waste, every useful program costs ten attempts. A typed program language
removes many stillborn moves by construction (\cref{sec:game}), and those
that remain usually fail before the costly fit. We therefore expect also-rans
to dominate, leaving $q_{\text{beat}}$, the factor only taste can raise.

\begin{figure}[t]
  \centering
  \includegraphics[width=\linewidth]{figures/fig4_price.pdf}
  \caption{\textbf{Most moves are wasted, and \TBD{the cost-per-win optimum
  sits at \ldots}.} (a)~Survival funnel: the share of moves from \TBD{n}
  pilot games on development families that run, score, beat the incumbent
  and survive re-verification, and the compute spent on moves that stop at
  each gate. (b)~Cost per win (log color scale) by writer size and reasoning
  effort; outlined, the measured minimum; hatched, no win. Error bars: 95\%
  bootstrap intervals over problems, \TBD{k} seeds. Drawn by
  \texttt{figures/fig4\_price.py} from logged pilot data.}
  \label{fig:price}
\end{figure}

\subsection{The correction tax and the reasoning dial}\label{sub:price-tax}

Small models are cheap per token and can be dear per win. Compute per token
falls about 190-fold from a dense 1.5T writer to an 8B one, so the small
model is cheaper per win unless its win rate per attempt is more than 190
times lower, repairs counted. That threshold flatters the small model:
repair turns re-read a growing context, and a writer that needs three times
the tokens per attempt breaks even at about 60, not 190.

Reasoning is the second dial. Reasoning tokens may cut stillborn moves, but
each is billed as output. \Cref{fig:price}b maps cost per win over writer
size and reasoning effort. A U-shaped curve is the hypothesis; the
measurement decides where, and whether, the minimum sits.

\subsection{Rewarding useful hallucination}\label{sub:price-hallucination}

A pretrained writer pulls toward the textbook answer, which is history's
taste. On a board whose incumbent is that answer, it is an also-ran by
definition. A \emph{useful hallucination} is a move the pretrained prior rates
unlikely and the referee confirms (Box~3). Four counterweights push the
writer toward such moves: a novelty bonus paid only inside the validity gate
(\cref{sec:game}), temperature and entropy schedules, a curriculum that
starts past known solutions, and a quality-diversity archive
\citep{lehman2011novelty,mouret2015mapelites}. \Cref{sec:method} states which
of these \EZ{} uses. Novelty also raises waste, so exploration gets a stated
share of games, \TBD{X\%}.

\begin{explainer}{Box 3: Why a referee makes hallucination safe}
  \label{box:hallucination}
  In text, a hallucination is an error nobody checks. In a game, an
  unfamiliar move is a conjecture, and the referee checks every one.
  Unfamiliar and verified is a discovery; unfamiliar and unchecked is
  confabulation. In March 2016, AlphaGo's own estimate put the chance that a
  human would play its Move~37 against Lee Sedol at about one in ten
  thousand \citep{metz2016twomoves}, and AlphaGo won the game. A referee does
  not make strange moves good. It makes them cheap to try, and it is only as
  trustworthy as what it measures (\cref{sec:method}).
\end{explainer}

\subsection{How many games?}\label{sub:price-games}

Three quantities differ by orders of magnitude.
\textbf{Tries to find one improvement:} $1/q$ on average (\cref{eq:taste}).
If tries were independent, $q = 1\%$ would need 69 tries for even odds and
299 for 95\% confidence; tries from one writer are correlated, so these are
guides, not guarantees.
\textbf{Experience to train a transferable player:} AlphaGo Zero generated
4.9 million self-play games in about three days \citep{silver2017mastering},
and AlphaZero played 44 million games of chess in about nine hours
\citep{silver2018general}. Our planning budget of 10,000
problems, 100 rounds and 8 attempts gives 8 million attempts, the same order.
It counts work, not learning: problem families matter more, and pilot
learning curves decide, read only within their measured range.
\textbf{Test problems to prove the claim:} on 150 independent problems, a win
rate near 60\% carries a 95\% interval of about $\pm 8$ points. We count
families, not datasets.

\subsection{A run, priced}\label{sub:price-run}

\Cref{tab:price} prices an 8-million-attempt run for six writers: about
\$12,000 with a 30B-class writer and about \$570,000 with a dense 1.5T one.
The dense 1.5T rollouts take $1.2\times10^{23}$~FLOPs, about 1\% of a $10^{25}$-FLOP
pretraining run, a scale more than 30 public models had reached by mid-2025
\citep{rahman2025over30,epoch2024data}. Mixture-of-experts (MoE) writers pay
compute on their active parameters, such as the 32B of 1T in Kimi~K2 \citep{kimiteam2025k2}
or the 37B of 671B in DeepSeek-V3 \citep{deepseekv3}. They need memory for
all of them: 1T one-byte weights fill at least 13 H100s.

\begin{table}[t]
  \centering
  \caption{\textbf{A full self-play run costs about \$12,000 with a 30B-class
  writer and about \$570,000 with a dense 1.5T writer; how many wins it buys
  depends on waste.} Projections before pilots; [low-cost, high-cost] in
  brackets. Base: 8 million attempts of 5,000 generated tokens; \$2.92 per
  H100-hour on 29 October 2026 \citep{ornn2026ocpi} [\$2.50, \$4.17;
  \citealp{thunder2026gpu}]; 989.5 dense BF16 TFLOP/s per GPU
  \citep{nvidia2026h100}; generation at 20\% of peak [30\%, 10\%], or 10\%
  [20\%, 7\%] for memory-bound mixture-of-experts (MoE) writers; updates on
  10\% of tokens at 40\% of peak [50\%, 30\%]. The 1.5T rows are scenarios.
  Excludes input tokens, repairs and the referee. \$k: thousands of dollars;
  source \texttt{scripts/cost\_model.py}.}
  \label{tab:price}
  \small
  \begin{tabular}{@{}l l S[table-format=2.2] S[table-format=2.2] l l@{}}
    \toprule
    Program writer & {Active} & {GPU-h per} & {\$ per} & {Rollouts, 8M} & {Updates,} \\
                   & {params} & {1,000 attempts} & {1,000 attempts} & {attempts (\$k)} & {full run (\$k)} \\
    \midrule
    Small dense & 8B & 0.11 & 0.33 & 2.6 [1.5, 7.5] & 0.4 [0.3, 0.7] \\
    Mid dense & 32B & 0.45 & 1.31 & 10.5 [6.0, 30.0] & 1.6 [1.1, 3.0] \\
    Large dense & 70B & 0.98 & 2.87 & 23.0 [13.1, 65.6] & 3.4 [2.4, 6.6] \\
    1T MoE & 32B & 0.90 & 2.62 & 21.0 [9.0, 42.8] & 1.6 [1.1, 3.0] \\
    1.5T MoE (scenario) & 50B & 1.40 & 4.10 & 32.8 [14.0, 66.9] & 2.5 [1.7, 4.7] \\
    1.5T dense (scenario) & 1.5T & 21.05 & 61 & 492 [281, 1405] & 73.8 [50.5, 140] \\
    \bottomrule
  \end{tabular}
\end{table}

The referee runs on CPUs, at 10 to 100 CPU-seconds per move and
\$0.015--\$0.045 per vCPU-hour \citep{doit2026c7i}. For a full run that is
\$330--\$990 at 10 seconds and \$3,300--\$9,900 at 100, enough to exceed an
8B writer's whole bill. List prices sit above our floor. On 30 October 2026,
Together AI listed \$0.60 per million output tokens for gpt-oss-120B (5.1B
active) and \$0.14 for an 8B Llama model
\citep{together2026pricing,openai2025gptoss}: about 7 and 2 times our floor
for the same model. The running example's single game cost \TBD{\$c}:
\TBD{t} generated tokens and \TBD{s} CPU-seconds of referee time.

\begin{keyidea}
  The cheapest way to make invention affordable is not cheaper tokens.
  It is fewer wasted moves.
\end{keyidea}

Every row of \cref{tab:price} scales with $1/q$ as surely as with
parameters: doubling $q$ halves the cost of a win at any writer size. The
method's job is now clear: raise $q$, so fewer moves are wasted.
